% TOP_PROBLEM: {"id": 30007267, "title": "The joint Dickman law for consecutive integers", "category_id": 1, "category": {"id": 1, "name": "number_theory", "display_name": "Number Theory", "description": "Properties of integers, prime numbers, Diophantine equations.", "slug": "number-theory", "order_index": 1, "created_at": "2026-07-31T15:26:25.670Z"}, "status": "open"}
% FIRST_POSTED: 2026-10-06
% ENTRY_KIND: statement_only
% !TeX program = lualatex
\documentclass[11pt]{article}
\usepackage[margin=25mm]{geometry}
\usepackage{fontspec}
\setmainfont{Latin Modern Roman}
\usepackage{amsmath,amssymb,mathtools}
\usepackage{unicode-math}
\setmathfont{Latin Modern Math}
\usepackage{xurl,hyperref}
\hypersetup{colorlinks=true,urlcolor=blue}
\setcounter{secnumdepth}{0}
\setlength{\parindent}{0pt}
\setlength{\parskip}{5pt}
\setlength{\emergencystretch}{3em}
\newcommand{\R}{\mathbb R}
\newcommand{\N}{\mathbb N}
\newcommand{\Z}{\mathbb Z}
\newcommand{\Q}{\mathbb Q}
\newcommand{\C}{\mathbb C}
\newcommand{\D}{\mathbb D}
\newcommand{\F}{\mathbb F}
\newcommand{\sref}[2]{\hyperlink{source:#1:#2}{[S#2]}}
\newcommand{\cataloguescope}[1]{#1}
\begin{document}
\section*{The joint Dickman law for consecutive integers}\label{30007267}
\subsection{Definitions and mathematical statement}
\paragraph{Statement recorded in the supplied source.}
For fixed a,b in (0,1), is the natural density of integers $n>=2$ with largest prime factor of n at most $n^a$ and that of n+1 at most $n^b$ equal to rho(1/a)rho(1/b)? Here $\rho (u)=1$ on [0,1] and u $\rho '(u)=-\rho (u-1)$ for $u>1$.
\subsection{Short English statement}
For fixed a,b in (0,1), is the natural density of integers $n>=2$ with largest prime factor of n at most $n^a$ and that of n+1 at most $n^b$ equal to rho(1/a)rho(1/b)? Here $\rho (u)=1$ on [0,1] and u $\rho '(u)=-\rho (u-1)$ for $u>1$.
\subsection{Sources}
\begin{itemize}
\item[S1] OpenAI, manuscript citation page. Source of the formulation motivating this local problem.
\url{https://github.com/openai/math/blob/adc7f1241b42e322a6451854ab7e4b4c146bf78a/preprints/The-joint-Dickman-law-for-consecutive-integers-September-24-2026/README.md}.
\end{itemize}
\end{document}
